% !TeX root = ../动手学强化学习笔记.tex
% 由主文档通过 \input 载入；请编译主文档。
\chapter*{课程结构与阅读指南}
\addcontentsline{toc}{chapter}{课程结构与阅读指南}

本书面向初次系统学习强化学习的读者，以“怎样通过交互学会作出长期有利的决策”为主线，依次讨论问题建模、价值估计、策略改进、函数近似与实验验证。学习起点是有限状态、有限动作的简单环境；神经网络在表格方法的基本逻辑建立之后引入。全书区分问题的数学模型、产生数据的环境程序、利用数据的学习算法以及最终执行动作的策略，避免把这些对象混为一谈。

\textbf{内容范围。}全书以本科读者能够复核的定义、递推、证明和数值例题建立主线。压缩映射、贝尔曼最优性、策略改进、策略梯度与奖励塑形等关键结论给出推导；随机逼近的深层收敛结果注明条件并作为引用使用。前三部分构成从表格方法到函数近似与深度方法的连续主线，第四部分提供进阶专题选读与贯穿各阶段的实验方法。深度算法强调目标、数据与实现之间的对应关系，进阶专题以基本模型、代表机制和适用边界为限，不要求穷尽各方向的算法与理论。各章习题与综合项目用于检验理解，附录提供部分计算题答案及证明提示。

\section*{先修知识与学习成果}

读者应具备一元与多元微积分、矩阵运算、离散概率、条件概率和数学期望的基础。条件期望、随机逼近与梯度估计等必要工具在第~\ref{ch:rl-math} 章按用途补充；开始表格方法时不要求先掌握深度学习。进入第~\ref{ch:rl-approximation} 章之后，需要理解参数化函数、损失函数、链式法则和小批量梯度更新。编程练习从数组、随机采样与循环开始，再过渡到自动微分。

\textbf{表格方法基础路径的学习成果。}读者应能把一个任务写成明确的序贯决策模型，推导贝尔曼递推与策略改进的关键关系，手算小规模价值函数，独立描述并实现动态规划、蒙特卡洛和时序差分方法；进一步解释多步回报、重要性采样与 Dyna 如何组织数据和更新。完成相应实验后，应能用已知真值检查价值估计，区分训练行为与最终策略的表现，并说明覆盖、步长和任务终止等条件对结论的影响。

\textbf{函数近似与深度方法主线的学习成果。}读者应能解释共享参数、自举与离策略采样带来的稳定性问题，推导策略梯度与优势估计的关键关系，辨认 DQN、Actor--Critic、PPO 及连续控制方法的目标、数据来源和梯度更新对象。实践上选择 DQN、PPO 或 SAC 中的一种完成实现与机制消融，能够分别检查目标值、梯度和回合边界，并报告独立评价结果及跨训练运行的波动；不以实现所有算法作为完成本阶段的必要条件。

\textbf{进阶专题选读的学习成果。}读者应能针对所选专题说明基础模型的哪项假设或决策条件发生了变化，写清新增的信息、目标或约束，解释一种代表机制及其适用边界，并用小规模例子或实验检验关键判断。选读成果以完成一个专题的建模、方法比较与证据分析为准，不要求同时掌握全部专题。

\section*{全书课程大纲}

\begin{enumerate}
  \item \textbf{第一部分：问题建模与数学基础。}第~\ref{ch:rl-introduction} 章“绪论”建立交互、反馈与长期目标的直观认识；第~\ref{ch:rl-math} 章“强化学习所需的数学基础”补充条件期望、矩阵和优化工具；第~\ref{ch:rl-bandits} 章“多臂赌博机与探索利用”研究最简单的评价性反馈；第~\ref{ch:rl-mdp} 章“马尔可夫决策过程”明确状态、转移与任务类型；第~\ref{ch:rl-bellman} 章“价值函数与贝尔曼方程”建立全书的递推语言。
  \item \textbf{第二部分：表格型强化学习。}第~\ref{ch:rl-dp} 章“动态规划”研究模型已知时的求解；第~\ref{ch:rl-mc} 章“蒙特卡洛方法”从完整回合估计价值；第~\ref{ch:rl-td} 章“时序差分学习”引入逐步自举；第~\ref{ch:rl-multistep} 章“多步学习与资格迹”统一不同长度的信用分配；第~\ref{ch:rl-offpolicy} 章“离策略学习与重要性采样”处理采样策略和目标策略不同的问题；第~\ref{ch:rl-planning} 章“模型学习与规划”把真实经验和模型模拟结合起来。
  \item \textbf{第三部分：函数近似与深度强化学习。}第~\ref{ch:rl-approximation} 章“价值函数近似”讨论泛化与稳定性；第~\ref{ch:rl-dqn} 章“深度价值方法”围绕深度 Q 网络组织算法；第~\ref{ch:rl-policy-gradient} 章“策略梯度方法”直接优化策略；第~\ref{ch:rl-actor-critic} 章“演员--评论家与优势估计”连接价值与策略；第~\ref{ch:rl-ppo} 章“信赖域方法与近端策略优化”讨论策略更新的尺度；第~\ref{ch:rl-continuous} 章“连续控制与最大熵强化学习”介绍 DDPG、TD3 与 SAC。
  \item \textbf{第四部分：进阶专题与综合实践。}第~\ref{ch:rl-deep-model} 章“深度模型学习与模型预测控制”研究近似模型如何支持规划及其误差；第~\ref{ch:rl-exploration} 章“探索、奖励设计与约束决策”分别讨论数据获取、奖励目标与策略可行性；第~\ref{ch:rl-pomdp} 章“部分可观测决策”处理隐状态、信息更新与记忆策略；第~\ref{ch:rl-hierarchy} 章“分层强化学习与时间抽象”讨论跨时间步的技能及其组合；第~\ref{ch:rl-offline} 章“模仿学习与离线强化学习”区分示范学习、固定数据下的策略优化，以及从示范或偏好推断奖励；第~\ref{ch:rl-multiagent} 章“多智能体强化学习”引入相互影响的决策者与不同的评价标准。上述六章按需选读，第~\ref{ch:rl-experiments} 章“实验方法、综合项目与论文阅读”则为各阶段共同使用的实践支撑。
\end{enumerate}

\section*{阅读路径与章节依赖}

\textbf{表格方法基础路径}覆盖第一、二部分，即第~\ref{ch:rl-introduction} 章至第~\ref{ch:rl-planning} 章。先以第~\ref{ch:rl-td} 章为阶段检查点，掌握建模、贝尔曼方程、动态规划、蒙特卡洛和一步时序差分，再学习多步方法、离策略校正与模型规划。同步阅读第~\ref{ch:rl-experiments} 章的环境验证与统计报告要求，完成项目一，并随着相应算法的学习逐步完成项目二；其中 Dyna 的对照在模型规划章之后进行。

\textbf{函数近似与深度方法主线}在基础路径之上覆盖第三部分，即第~\ref{ch:rl-approximation} 章至第~\ref{ch:rl-continuous} 章。函数近似与离策略学习支撑深度价值方法；多步回报支撑优势估计，策略梯度支撑演员--评论家与后续策略优化。阅读时先弄清各方法的目标与更新，再按任务动作空间选择 DQN、PPO 或 SAC 完成第~\ref{ch:rl-experiments} 章的项目三。前三部分的理论学习与这一实现项目共同构成完整主线，项目所选算法不替代对其他核心方法的概念理解。

\textbf{进阶专题选读}从第四部分前六章中按问题需要选择，不要求按章顺序全部学习，也不作为基础路径的前提。深度模型学习依赖模型规划、函数近似与优化；探索专题承接赌博机与价值学习，奖励塑形依赖回报和贝尔曼方程，约束决策还需理解约束优化与拉格朗日乘子；部分可观测决策依赖条件概率与 MDP，涉及循环网络时再补充序列训练知识；分层方法依赖多步回报与规划。模仿学习的行为克隆需要监督学习基础，离线策略学习还依赖离策略校正与函数近似，奖励学习部分则需结合奖励塑形理解可辨认性；多智能体部分需要博弈论基本概念，集中训练的评论家结构还需演员--评论家基础。完成相应先修后，可用所选专题的章末习题或第~\ref{ch:rl-experiments} 章的项目四开展专题验证。

若用于一个学期的教学，应先确定课程面向表格方法基础还是深度方法进阶，再据此分配推导、练习与实验时间。基础课程以第一、二部分和项目一、二为核心，进阶课程在已掌握基础路径的前提下以第三部分和项目三为核心；第四部分的六个专题按时间与目标选取。实验方法在两类课程中均为必学内容，具体技术检查随对应算法逐步引入。

贯穿全书的例子分为四组：伯努利赌博机用于观察探索与估计误差；两状态过程、随机游走和网格世界用于手算与验证；低维连续控制用于连接函数近似和策略优化；固定数据集与局部观测迷宫用于检验进阶方法的假设。每章习题按照概念辨析、数值计算、推导证明和实现分析组织，实验结果必须来自实际运行，不能用预期趋势代替观测。

\section*{符号与理论边界}

全书沿用 $S_t,A_t,R_{t+1}$ 的交互时序：先在状态 $S_t$ 选择动作 $A_t$，再获得奖励 $R_{t+1}$ 并转移到 $S_{t+1}$。小写字母表示具体取值；$v_\pi,q_\pi$ 表示策略的真实价值，$V,Q$ 表示学习得到的估计，$V_w,Q_w$ 表示参数为 $w$ 的近似。策略参数统一记为 $\theta$，价值函数参数记为 $w$；学习率记为 $\alpha_t$，折扣因子记为 $\gamma$，资格迹或优势估计中的衰减参数记为 $\lambda$。各章在首次使用专用记号时给出其定义，章节间沿用同一交互时序。

基础理论优先采用有限状态、有限动作、平稳动力学、奖励有界且 $0\leq\gamma<1$ 的折扣模型。有限时域、无折扣回合、平均奖励、连续空间和部分可观测模型分别说明额外条件。收敛到真实价值、收敛到最优策略、参数达到驻点以及有限样本中表现较好是不同结论；表格算法的保证不得直接移用于神经网络算法。基础阅读可参照 Sutton 与 Barto 的教材\cite{sutton-barto}，专题算法在对应章节给出原始文献。

